\hypertarget{functional-breadth-map-filter-and-reduce}{%
\chapter{\texorpdfstring{Functional Breadth: \texttt{map},
\texttt{filter}, and
\texttt{reduce}}{Functional Breadth: map, filter, and reduce}}\label{functional-breadth-map-filter-and-reduce}}

\hypertarget{the-operator-module-integration}{%
\subsection{\texorpdfstring{94.1 The \texttt{operator} Module
(Integration)}{94.1 The operator Module (Integration)}}\label{the-operator-module-integration}}

As seen in Chapter 34, combining \passthrough{\lstinline!map!} with the
\passthrough{\lstinline!operator!} module is often faster than lambdas.

\begin{lstlisting}[language=Python]
from operator import add
result = list(map(add, [1, 2, 3], [4, 5, 6])) # [5, 7, 9]
\end{lstlisting}

\hypertarget{reduce-and-accumulate}{%
\subsection{\texorpdfstring{94.2 \texttt{reduce} and
\texttt{accumulate}}{94.2 reduce and accumulate}}\label{reduce-and-accumulate}}

\begin{itemize}
\tightlist
\item
  \textbf{\passthrough{\lstinline!functools.reduce!}}: Collapses a
  sequence to a single value by applying a binary function cumulatively.
\item
  \textbf{\passthrough{\lstinline!itertools.accumulate!}}: Similar to
  reduce, but yields every intermediate result.
\end{itemize}

\begin{center}\rule{0.5\linewidth}{0.5pt}\end{center}

\hypertarget{phase-xxii-visualization-and-interface-engineering}{%
\section{Phase XXII: Visualization and Interface
Engineering}\label{phase-xxii-visualization-and-interface-engineering}}

\hypertarget{chapter-95-turtle-graphics-the-educational-engine}{%
\chapter{Chapter 95: Turtle Graphics: The Educational
Engine}\label{chapter-95-turtle-graphics-the-educational-engine}}

The \passthrough{\lstinline!turtle!} module is a built-in toolkit for
turtle graphics, providing an excellent way to visualize algorithms and
teach geometry.

\hypertarget{the-virtual-screen-and-the-turtle}{%
\subsection{95.1 The Virtual Screen and the
Turtle}\label{the-virtual-screen-and-the-turtle}}

\begin{itemize}
\tightlist
\item
  \textbf{The Turtle}: A stateful cursor that maintains a position, a
  heading, and a ``pen'' (up or down).
\item
  \textbf{The Screen}: A window where the turtle draws.
\end{itemize}

\hypertarget{recursive-fractals-with-turtle}{%
\subsection{95.2 Recursive Fractals with
Turtle}\label{recursive-fractals-with-turtle}}

Because the turtle's state is easily managed, it is perfect for drawing
recursive structures like the Koch Snowflake or the Sierpinski Triangle.

\begin{center}\rule{0.5\linewidth}{0.5pt}\end{center}
